% 教材查阅指南示例（微观经济学，只能带教材）。教材名和页码都是虚构的，只示范格式
\documentclass[a4paper,fontset=windows]{ctexart}
\ifdefined\SBSIZE\else\def\SBSIZE{10}\fi
\ifdefined\SBLEAD\else\def\SBLEAD{14}\fi
\ifdefined\SBMARGIN\else\def\SBMARGIN{18}\fi
\usepackage[margin=\SBMARGIN mm]{geometry}
\usepackage{amsmath,amssymb,enumitem,booktabs,longtable,xcolor}
\makeatletter
\renewcommand\normalsize{\@setfontsize\normalsize{\SBSIZE}{\SBLEAD}%
  \abovedisplayskip 4\p@ \@plus2\p@ \@minus2\p@
  \belowdisplayskip \abovedisplayskip
  \abovedisplayshortskip 2\p@ \@plus1\p@
  \belowdisplayshortskip 2\p@ \@plus1\p@
  \let\@listi\@listI}
\makeatother
\normalsize
\pagestyle{plain}
\setlength{\parindent}{0pt}
\setlength{\parskip}{3pt}
\setlist{nosep,leftmargin=1.4em}
\newlength{\sbw}
\newcommand{\kw}[1]{\ifdefined\SBFILL\settowidth{\sbw}{#1}\underline{\hspace*{\dimexpr\sbw+0.6em\relax}}\else\textbf{#1}\fi}

\begin{document}
\begin{center}
{\LARGE\bfseries 微观经济学　教材查阅指南}\\[4pt]
{\small 教材：《微观经济学原理》（示例）　考试：2027-01-08　规则：只能携带教材}
\end{center}

页码均为教材上印的页码。

\begin{longtable}{@{}p{0.26\linewidth}p{0.2\linewidth}p{0.4\linewidth}c@{}}
\toprule
考点 & 教材位置 & 在那里找什么 & 权重 \\
\midrule
\endhead
需求价格弹性的计算 & 第5章5.1节，p.92--98 & 公式（5.1）；例5.1（奶茶涨价）；表5.2 & 高 \\
弹性与总收益 & 第5章5.3节，p.103--107 & 图5.4（总收益随价格变化）；例5.3 & 高 \\
影响弹性的因素 & 第5章5.2节，p.99--102 & 三个因素的小标题；汽油短期与长期的例子 & 高 \\
均衡价格与数量 & 第4章4.3节，p.74--80 & 图4.5；例4.2（联立方程解均衡） & 高 \\
需求移动与沿曲线移动 & 第4章4.2节，p.66--73 & 图4.3和图4.4；本节末的对比框 & 高 \\
消费者剩余 & 第7章7.1节，p.130--136 & 图7.1；公式（7.2）；例7.1 & 中 \\
生产者剩余 & 第7章7.2节，p.137--140 & 图7.3 & 低 \\
\bottomrule
\end{longtable}

\section*{附录里的表（考试可能要查）}
\begin{itemize}
\item 常用公式汇总：附录A，p.412
\item 术语中英对照：附录B，p.418
\end{itemize}

\end{document}
